\subsection{模的张量代数的定义}

设 A 是一个交换环，M 是一个 A-模。对于每个整数 \(n \geqslant 0\) ，A-模，即 n 个都等于 M 的模的张量积（也称为 M 的第 n 次张量幂），记为 \(\bigotimes^{n} M\) ，或者 \(M^{\otimes n}\) ，或者 \(T^{n}(M)\) ，或者 \(T_{A}^{n}(M)\) ，或者 \(\operatorname{Tens}^{n}(M)\) ；于是 \(T^{1}(M) = M\) ；同时我们记 \(T^{0}(M) = A\) 。A-模，即直和 \(\bigoplus_{n \geqslant 0} T^{n}(M)\) ，记为 \(T(M)\) 或 \(\operatorname{Tens}(M)\) 。我们将在 \(T(M)\) 上定义一个类型为 N 的分次 A-代数结构，方法是：对每一对有序整数 \(p \geqslant 0\) , \(q \geqslant 0\) ，定义一个 A-线性映射

\[
m _ {p q} \colon \mathsf {T} ^ {p} (\mathbf {M}) \otimes_ {\mathbf {A}} \mathsf {T} ^ {q} (\mathbf {M}) \rightarrow \mathsf {T} ^ {p + q} (\mathbf {M})
\]

（§ 3，第 1 号，注记）。当 \(p > 0\) 且 \(q > 0\) 时， \(m_{pq}\) 是结合性同构（II，§ 3，第 9 号），而当 \(p = 0\) （相应地 \(q = 0\) ）时， \(m_{0,q}\) 是 \(A \otimes_A T^q(M)\) 到 \(T^q(M)\) 上的典范同构（相应地， \(m_{p,0}\) 是 \(T^p(M) \otimes_A A\) 到 \(T^p(M)\) 上的典范同构）（II，§ 3，第 4 号，命题 4）。然后，对于 \(x_t \in M\) , \(\alpha \in A\) ,

\[
\left\{ \begin{array}{r l} (x _ {1} \otimes \dots \otimes x _ {p}) \cdot (x _ {p + 1} \otimes \dots \otimes x _ {p + q}) & \\ = x _ {1} \otimes \dots \otimes x _ {p} \otimes x _ {p + 1} \otimes \dots \otimes x _ {p + q} \\ \alpha \cdot (x _ {1} \otimes \dots \otimes x _ {p}) = \alpha (x _ {1} \otimes \dots \otimes x _ {p}). \end{array} \right.\tag{1}
\]

显然，如此定义在 \(\mathsf{T}(\mathbf{M})\) 上的乘法满足结合律，并以单位元 \(1\) （即 \(\mathbf{A} = \mathsf{T}^0 (\mathbf{M})\) 的单位元）为单位元素。

定义1. 对于交换环A上的每个模M，M的张量代数，记为T(M)或Tens(M)，或 \(\mathsf{T}_{\mathbf{A}}(\mathbf{M})\) ，是代数 \(\bigoplus_{n\geqslant 0}\mathsf{T}^n (\mathbf{M})\) ，其乘法由(1)定义。典范单射 \(\phi :\mathsf{T}^{1}(\mathbf{M})\to \mathsf{T}(\mathbf{M})\) （II，§ 1，第 12 号）（也记作 \(\phi_{\mathbf{M}}\) ）称为M到T(M)的典范单射。

命题1. 设E是一个（含单位元的）A-代数，且 \(f:M\to E\) 为一个A-线性映射。存在唯一的一个A-代数同态 \(g:T(M)\to E\) ，使得 \(f=g\circ\phi\) .

换言之， \((\mathsf{T}(\mathbf{M}), \phi)\) 是泛映射问题（集合论，IV，§3，第1号）的一个解，其中 \(\Sigma\) 是A-代数结构的类， \(\alpha\) -映射是从模 M 到 A-代数的 A-线性映射。注意，这里完全不涉及 \(\mathsf{T}(\mathbf{M})\) 上的分次问题。

对于每个有限族 \((x_{i})_{1\leqslant i\leqslant n}\) ，它由 M 中的 n 个元素组成，根据 T(M) 中乘积的定义，有 \(x_{1}\otimes x_{2}\otimes\cdots\otimes x_{n}=\phi(x_{1})\phi(x_{2})\ldots\phi(x_{n})\) ；于是必然有 \(g(x_{1}\otimes x_{2}\otimes\cdots\otimes x_{n})=f(x_{1})f(x_{2})\ldots f(x_{n})\) （对 \(n\geqslant1\) ）且 \(g(\alpha)=\alpha e\) （若 e 是 E 的单位元）（对 \(\alpha\in A\) ），这就证明了 g 的唯一性。反之，注意到对所有 n\textgreater0，映射

\[
(x _ {1}, \dots , x _ {n}) \mapsto f (x _ {1}) f (x _ {2}) \dots f (x _ {n})
\]

从 \(\mathbf{M}^n\) 到 \(\mathbf{E}\) 的映射是 A-多重线性的；因此它对应着一个 A-线性映射 \(g_n: \mathsf{T}^n(\mathbf{M}) \to \mathbf{E}\) ，使得

\[
g \left(x _ {1} \otimes x _ {2} \otimes \dots \otimes x _ {n}\right) = f \left(x _ {1}\right) f \left(x _ {2}\right) \dots f \left(x _ {n}\right).\tag{2}
\]

我们还定义映射 \(g_0 \colon \mathsf{T}^0(\mathbf{M}) \to \mathbf{E}\) ，它等于 \(\eta_{\mathbf{E}}\) （§ 1，第 3 号），换言之， \(g_0(\alpha) = \alpha e\) （对 \(\alpha \in \mathbf{A}\) ）。设 \(g\) 是 \(\mathsf{T}(\mathbf{M})\) 到 \(\mathbf{E}\) 中的唯一的 A-线性映射，它在 \(\mathsf{T}^n(\mathbf{M})\) 上的限制为 \(g_n\) ( \(n \geqslant 0\) )；显然 \(g \circ \phi = g_1 = f\) ，剩下的只需验证 \(g\) 是 A-代数同态。由构造有 \(g(1) = e\) ，且由线性性只需证明 \(g(uv) = g(u)g(v)\) 对于 \(u \in \mathsf{T}^p(\mathbf{M})\) 和 \(v \in \mathsf{T}^q(\mathbf{M})\) ( \(p > 0, q > 0\) )成立；而现在由公式 (1) 和 (2) 可知，这一关系式当 \(u = x_1 \otimes x_2 \otimes \cdots \otimes x_p\) 且 \(v = x_{p+1} \otimes \cdots \otimes x_{p+q}\) （其中 \(x_i\) 属于 \(\mathbf{M}\) ）时成立。因此由线性性，它对于 \(u \in \mathsf{T}^p(\mathbf{M})\) 和 \(v \in \mathsf{T}^q(\mathbf{M})\) 成立。

注记。假设 \(\mathbf{E}\) 是类型为 \(\mathbf{Z}\) 的分次A-代数，其分次为 \((\mathbf{E}_n)\) ，并且还假设

\[
f (\mathbf {M}) \subset \mathrm{E} _ {1}.\tag{3}
\]

于是由(2)可知 \(g(\mathsf{T}^p (\mathbf{M})) \subset \mathrm{E}_p\) （对所有 \(p \geqslant 0\) ），因此 \(g\) 是分次代数同态。

